\subsection{Universal sections of polarization torsors}
\paragraph{Conventions.}
Throughout this subsection $2$ is invertible on the base.  We use the
fppf relative Picard sheaf, $B=A^\vee=\operatorname{Pic}^0_{A/S}$,
the usual Poincaré biduality $A=B^\vee$, and
\[
 \lambda_L(a)=[t_a^*L\otimes L^{-1}].
\]
Thus an ample $L$ gives a positive polarization $\lambda_L:A\to B$.
The torsor $T_\delta=\operatorname{Pic}^{\delta}_{A/S}$ is a torsor under
$B$ by tensor product.  A section here means a section of this sheaf
over the moduli stack, and hence a construction invariant under
isomorphisms and commuting with arbitrary base change.  It is not the
assertion that each individual algebraically closed fibre has a point.

\paragraph{G1 (Proved: the complete existence criterion).}
Let $D=(d_1\mid\cdots\mid d_n)$ be a polarization type, and put
$r=\#\{i: d_i\text{ is odd}\}$.  The universal $T_\delta$ over the
moduli stack of type $D$ over $\mathbb Z[1/2]$ has a section if and only
if $r\leq1$.  The construction in the affirmative cases is canonical.
It requires no inversion of the odd divisors of $\prod_i d_i$.
More intrinsically, its hypothesis is that the alternating pairing
on $A[2]$ induced by $\lambda$ has rank at most two.  This formulation
also specifies the assertion on integral models of the polarization
moduli problem at primes dividing its degree: it depends only on the
locally constant type of the étale group $A[2]$, not on terminology for
the non-étale part of the polarization kernel.

We first recall precisely the input about symmetric line bundles.
The symmetric classes in $T_\delta$ form an étale $B[2]$-torsor.  The
normalized symmetry of a line bundle identifies this torsor with the
torsor of quadratic refinements of
\[
 b_\lambda(x,y)=e_2(x,\lambda(y)),\qquad x,y\in A[2].
\]
We write $\mu_2$ additively as $\mathbb F_2$, so the refinement identity
is $q(x+y)=q(x)+q(y)+b_\lambda(x,y)$.  This is the scheme-valued
statement of Poonen--Rains, Proposition 3.6 and diagram (3.7),
pp.~1312--1313 \cite{DGJ_PoonenRains2011}; the exact sequence preceding
it identifies its fibre over $\lambda$ with the stated torsor.
In particular this is an identification of sheaves over arbitrary
bases, rather than a bijection only on geometric points.

Let $R$ be the radical of $b_\lambda$.  It is a finite étale subgroup
of $A[2]$, because the rank of this map of finite étale groups is locally
constant.  If $r=0$, take $q=0$.  If $r=1$, define
\[
 q(x)=\begin{cases}0&x\in R,\\1&x\notin R.\end{cases}
\]
This is a morphism of finite étale schemes: both loci are open and
closed.  It is a quadratic refinement, as can be checked on the
two-dimensional quotient $A[2]/R$.  In that quotient distinct nonzero
vectors have pairing one, and their sum is the third nonzero vector.
The other cases in the refinement identity are immediate.  The
definition uses no choices and therefore gives the asserted section
of the symmetric Picard torsor, and of $T_\delta$.

For necessity we may pull a proposed section back to characteristic
zero.  The automorphism $[-1]$ of every polarized abelian scheme
forces the proposed section to be symmetric.  Over the Siegel space
of marked complex polarized abelian varieties, the resulting section
of a finite étale torsor is constant.  Descent would therefore give
a quadratic refinement fixed by the integral group preserving the
alternating lattice with basis $e_i,f_i$ and
$E(e_i,f_j)=d_i\delta_{ij}$.  This description of the complex moduli
stack as the quotient of Siegel space by this group is recorded, with
the definition of the group, in \cite[\S\S5.3--5.4, pp.~30--33]{DGJ_Laga2025}.

For every integral lattice vector $v$, the integral transvection
$x\mapsto x+E(x,v)v$ preserves $E$ and has integral inverse
$x\mapsto x-E(x,v)v$.  In the subspace $W$ of $A[2]$ spanned by the
pairs with odd $d_i$, its reductions include the transvection along
every $v\in W\setminus\{0\}$.  If $q$ were invariant, choose $x\in W$
with $b_\lambda(x,v)=1$.  Invariance and the quadratic identity give
$q(x)=q(x+v)=q(x)+q(v)+1$, hence $q(v)=1$ for every nonzero $v\in W$.
When $r\geq2$, choose independent orthogonal vectors $v,w\in W$.
Then $q(v+w)=q(v)+q(w)=0$, a contradiction.  Thus no section exists
in these cases, already over $\mathbb C$.

The determinant comparisons below use the following uniqueness fact.  There is
no nonzero natural section of the universal $A[2]$ or $A^\vee[2]$ over
the moduli problem of polarized abelian schemes of given dimension
and polarization degree, or over its locus where the rank of
$b_\lambda$ is fixed.  In characteristic zero, the integral lattice
group contains an $\operatorname{SL}_2(\mathbb Z)$ on each pair
$e_i,f_i$, even when $d_i$ is even.  Its action on the corresponding
two-dimensional summand modulo two has no fixed nonzero vector.
There are therefore no invariant vectors in the whole lattice modulo
two, or in its dual.  This proves the assertion on each complex
component.  It implies the characteristic-zero assertion over any
field by spreading an object to a finitely generated field and
embedding that field in $\mathbb C$.
For a polarized abelian variety in positive characteristic we use
the polarized lifting theorem of Mumford--Norman--Oort, stated as
Theorem (8.3).1 in \cite[p.~29]{DGJ_Oort2014}: it lifts over an integral
mixed-characteristic ring.  The degree and the rank on two-torsion
are constant on this ring.  Two sections of a finite étale separated
scheme that agree generically agree everywhere on this connected
base, since their equalizer is open and closed.  The asserted section
is thus zero also on the special fibre.  Finally equality of sections
of a finite étale scheme can be tested on all geometric points, so
this proves the assertion over arbitrary, including nonreduced,
bases.  Thus two constructions of sections of $T_\delta$ which are
natural on this entire degree locus (or the indicated rank locus)
are equal: their difference is such a two-torsion section because
each is fixed by $[-1]^*$.  This is the form of uniqueness needed
below; no claim about extra vertical components in an unspecified
integral model of a fixed-type moduli stack is used.

The finite computation accompanying this paper enumerates refinements
and transvection invariants in ranks zero, two, four and six, with
radicals of dimensions zero or two.  It confirms the small cases;
the argument above proves every rank and does not rely on extrapolation.

\subsection{The determinant construction and its signs}
\paragraph{G2a (Proved: R9's boxed relation, with typed signs).}
Let $\mathcal P$ be the normalized Poincaré bundle on $A\times_S B$.
For a relatively ample line bundle $L$ of class $\delta$, write
\[
 \Phi(L)=R\operatorname{pr}_{B*}
       (\mathcal P\otimes\operatorname{pr}_A^*L),\qquad
 F_\delta(L)=\det\Phi(L)^{-1},\qquad \chi=\chi(L).
\]
These are relative constructions.  The complex is perfect; in the
ample case it is a vector bundle of rank $\chi$.  Tensoring $L$ by a
bundle from $S$ changes its determinant by a bundle from $S$ only,
so $F_\delta$ defines a morphism of relative Picard sheaves.
The theorem of the square gives, up to bundles from the base,
\[
 (1,\lambda_L)^*\mathcal P
   =m^*L\otimes\operatorname{pr}_1^*L^{-1}
          \otimes\operatorname{pr}_2^*L^{-1}.
\]
Changing variables $(x,a)\mapsto(x+a,a)$ therefore gives
$\lambda_L^*\Phi(L)=\pi_*L\otimes L^{-1}$ up to the same harmless
base factors.  Taking determinants shows
$\lambda_L^*F_\delta(L)=L^\chi$ in the relative Picard sheaf.
Consequently, if $\delta'$ is the class of $F_\delta(L)$, then
\[
 \lambda_\delta\lambda_{\delta'}=[\chi]_B,
 \qquad
 \lambda_{\delta'}\lambda_\delta=[\chi]_A.
\]
Indeed the formula for a pulled-back line bundle first gives
$\lambda_\delta\lambda_{\delta'}\lambda_\delta
=\chi\lambda_\delta$; the isogeny $\lambda_\delta$ is faithfully flat
and can be cancelled on the right.  Duality gives the other equality.
Thus $\delta'$ is positive, whereas the determinant itself has class
$\varphi(\delta)=-\delta'$.  This proves precisely the negative sign
in the box of 59:10.

The resulting torsor morphism and its equivariance are
\[
 F_\delta:T_\delta\longrightarrow T_{\delta'},\qquad
 F_\delta(L\otimes\beta)=F_\delta(L)\otimes\lambda_{\delta'}(\beta),
 \quad \beta\in B.
\]
This follows from
$\Phi(L\otimes\beta)=t_\beta^*\Phi(L)$.
In particular, with our stated translation convention the linear
part is $+\lambda_{\delta'}$.  The page's use of $T_{-\delta}$ can be
typed without an ambiguous suppression of signs: the canonical
$B$-equivariant isomorphism
\[
 T_\delta\longrightarrow T_{-\delta},\qquad
 L\longmapsto([-1]^*L)^{-1}
\]
has linear part $+1$, whereas $L\mapsto L^{-1}$ has linear part $-1$.
Using the first gives the displayed version of (*) on 59:12.

The equalities also give
$\chi(\delta')=\chi^{n-1}$ and
$\delta''=\chi^{n-2}\delta$ for $n\geq2$, as equalities of
Néron--Severi sections.  To see the first, take degrees of the
isogenies and use positivity to select the positive square root.
The second follows after applying the same relation to $B$ and
using the injection of the Néron--Severi sheaf into symmetric
homomorphisms.  In dimension one the double transform is the
canonical degree-one polarization, so the expression with a negative
exponent is interpreted in this sense, not as tensoring a torsor
by a nonintegral integer.

\paragraph{G2b (Proved: the ``à vérifier'' comparisons).}
For every Néron--Severi section $\eta$, the norm
\[
 c_\eta=L\otimes[-1]^*L\in T_{2\eta}
\]
is independent of the local choice of $L\in T_\eta$: a change by
$\beta\in A^\vee$ cancels with $[-1]^*\beta=\beta^{-1}$.
It is a canonical section, compatible with arbitrary base change.
The determinant commutes with $[-1]^*$, as is seen by pulling back
the Poincaré bundle and changing variables on $A$.  It therefore
sends symmetric classes to symmetric classes.  Choose such a class
étale locally, which is possible because $2$ is invertible.  Then
$2L=c_\delta$ and $2F_\delta(L)=c_{\delta'}$.  Thus the isomorphism
of squared torsors induced by $F_\delta$ carries the canonical
section to the canonical section.  This proves the margin question
at 59:12, including its formulation as an isomorphism of the
corresponding two-torsion torsors.

We spell out why the additional determinant trivializations on
59:13--14 have the same normalization, rather than inferring it from
an equality of torsor classes.  Determinants of perfect complexes,
tensor powers, pullbacks, and their compositions are all natural
constructions on the full polarized moduli problem.  The comparison
of $\delta$ with $-\delta$ used there produces a section of
$T_{2\delta'}$; relative Serre duality supplies the comparison also
for the antiample class.  Namely, up to a base line,
\[
 F_{-\delta}(L^{-1})
       =\bigl([-1]^*F_\delta(L)\bigr)^{(-1)^{n-1}}.
\]
Both this comparison section and $c_{\delta'}$ are invariant under
$[-1]$; their difference is consequently a natural section of
$A[2]$.  The uniqueness argument in G1 makes it zero.  This proves
the margin question on the canonical square at 59:13.

For the double transform, put $a=\chi^{n-2}$ with $n\geq2$.
The composite $F_{\delta'}F_\delta$ has target $T_{a\delta}$ and
linear part $[a\chi]_B$.  Comparing it with the tensor-power map
$L\mapsto L^a$ therefore supplies a natural trivialization of
$T_{a(\chi-1)\delta}$, with reversal of the entire sign if the
opposite comparison is chosen.  The preceding square construction
also trivializes $T_{2a\delta}$.  If $\chi$ is even, $\chi-1$ is odd;
combining these trivializations with integral Bézout coefficients
gives a natural section of $T_{a\delta}$.  It is invariant under
$[-1]$, as is every step of the construction.  When $n\geq3$, $a$
is even and its comparison with $c_{(a/2)\delta}$ again differs by
a universal two-torsion section; hence they agree.  When $n=2$,
$a=1$ and $\chi$ even forces at most one odd elementary divisor.
The resulting section therefore agrees with the section of G1,
by the same uniqueness lemma on the full even-degree surface locus.
Changing Bézout coefficients gives the same result by this
uniqueness.  These statements establish the determinant
trivializations as actual sections, not merely vanishing classes
in $H^1(S,B)$.

There are two literal slips to retain in the record.  The word
``ppcm'' on 59:13--14 must be ``pgcd'' for the Bézout deduction:
for $n\geq2$ every $m^n-m^{n-1}$ is even and the value at $m=-1$
is $\pm2$.  Also the assertion on 59:13 that an even $\chi(\delta)$
implies divisibility of $\delta$ by two is false.  On $E\times E$,
the exterior product of degree-one and degree-two line bundles has
type $(1,2)$ and Euler characteristic two, but its restriction to
the first factor has odd degree, so its class is not twice another
class.  G1 supplies exactly the needed correction in dimension two.
The $B$ subscripts for the torsors of class $\delta$ on 59:14 must
be $A$; the class $\delta$ was defined on $A$.

\subsection{The Albanese sign and the canonical theta class}
\paragraph{G3 (Proved).}
At 59:23 the selected convention is
$\delta(x,y)=[N_yN_x^{-1}]$, followed by
``Nous prendrons la première convention''.  For a smooth proper
geometrically connected curve $C/S$, let $J=\operatorname{Pic}^0_{C/S}$
and let $\lambda_\Theta:J\to J^\vee$ be its positive theta
polarization with our translation convention.  Then the map of
59:24 is
\[
 c:J\longrightarrow J^\vee,\qquad c=-\lambda_\Theta.
\]
Thus the correction on that page, ``Non, c'est sans doute la
symétrique de la $\Theta$-polarisation'', has the right sign.
To verify it, choose a section $x$ locally and let
$a:C\to J$, $y\mapsto[y-x]$.  The definition by the Poincaré
bundle identifies $c$ with the inverse of
$a^*:J^\vee\to J$.  Jacobian autoduality gives
$a^*\lambda_\Theta=-1$; see Milne, Lemma 6.9 and
Remark 6.10(c), p.~23 of the revised notes
\cite{DGJ_MilneJac}.  We inspected the rendered page because text
extraction drops the minus sign in this lemma.  The equality over
families follows from its fibrewise equality and rigidity of
homomorphisms of abelian schemes.  This includes nonreduced bases;
the relative Hom functor is unramified.  In particular transporting
$\operatorname{Alb}^1_{C/S}$ from $J^\vee$ to $J$ by
$\lambda_\Theta^{-1}$ gives $\operatorname{Pic}^{-1}_{C/S}$, in
agreement with the sign warning on 59:9.

For a principally polarized abelian scheme $(A,\delta)$ there is a
canonical principally polarized torsor $(P,\Theta)$ as described
on 59:26--27.  One construction is $P=T_\delta$, with its structural
group transported from $A^\vee$ to $A$ by $\lambda_\delta$.
The universal line bundle has a rank-one pushforward, and its
evaluation at the zero section defines the theta divisor on $P$.
For any other such torsor, sending a point $p$ to the line bundle
pulled back from $\Theta$ by $a\mapsto a+p$ gives the required
isomorphism with $T_\delta$; its linear part is
$\lambda_\delta$, hence it is an isomorphism.  Its uniqueness
follows since a translation preserving $\Theta$ lies in the kernel
of this principal polarization.  These arguments are fppf local
and hence descend to arbitrary bases.

The involution $L\mapsto[-1]^*L$ of $T_\delta$ preserves this
theta divisor: it preserves the condition that the unique section
vanish at zero.  It has linear part $-1$ on its structural group.
By the just-proved uniqueness, it is precisely the theta involution
used on 59:28 to define $K\in P^{\otimes2}$.  Consequently
\[
 K=L+[-1]^*L=c_\delta\in T_{2\delta}.
\]
This proves that page's normalization question directly.

For a curve the theta torsor is
$\operatorname{Pic}^{g-1}_{C/S}$ with its effective theta divisor.
Relative Serre duality identifies its theta involution as
\[
 L\longmapsto\omega_{C/S}\otimes L^{-1}.
\]
The equality holds scheme-theoretically, since the theta divisor
is the determinant divisor of cohomology and Serre duality
identifies the corresponding determinants.  Uniqueness of the
theta involution now gives
$K=[\omega_{C/S}]$ in $\operatorname{Pic}^{2g-2}_{C/S}$.
No assertion about generators of the Picard group of the moduli
stack of curves is needed.  The opposite structural-group
identifications mentioned on 59:27 are related by this very
involution; the minus sign of $c$ must not be erased when writing
an equivariant map.
